% Unit 135 terjemahan Bahasa Indonesia dari appendix1.tex baris 1--120.
% Sumber: Methods of Algebra, Volume 2, commit
% 9a5803ff77dd3257484cb177f851a73770a59dd3 (CC BY 4.0).
% stable-unit-id: o014.aljabr2.appendix1.yoneda-density
% Cakupan: pembukaan Lampiran A dan Bagian A.1 lengkap.

% segment-id: o014.li2.appendix1.yoneda-density.h001
\chapter{Materi lanjutan tentang kategori abelian}\label{sec:app-Abel}
\hypertarget{o014.aljabr2.appendix1.yoneda-density}{ }
% segment-id: o014.li2.appendix1.yoneda-density.p001
Lampiran ini merupakan kumpulan materi yang berkaitan langsung maupun tidak
langsung dengan kategori abelian, dan sifatnya hampir seperti ``gado-gado''.
Pengetahuan prasyaratnya pada prinsipnya tidak melampaui bagian pertama
\CHref{sec:categories} dan sebagian kecil \CHref{sec:Abel-cat}.

% segment-id: o014.li2.appendix1.yoneda-density.p002
Bagian~\ref{sec:Yoneda} mula-mula meninjau singkat pembenaman Yoneda.
Teorema Kepadatan~\ref{prop:Yoneda-density} di dalamnya merupakan fakta
standar teori kategori, tetapi sering diabaikan dalam buku ajar. Alat-alat
ini akan sering dipakai dalam \S\ref{sec:ind-pro}.

% segment-id: o014.li2.appendix1.yoneda-density.p003
Walaupun objek kompak dan kategori tersajikan yang diperkenalkan dalam
\S\ref{sec:cpt-objects} tidak mutlak diperlukan di bagian lain buku ini,
keduanya menyediakan bahasa yang umum dan praktis. Pembahasan kardinal
reguler di sana juga membuka jalan bagi \S\ref{sec:Grothendieck-cat}.

% segment-id: o014.li2.appendix1.yoneda-density.p004
Dua bagian berikutnya berkaitan langsung dengan kategori abelian. Teorema
Gabriel--Popescu dalam \S\ref{sec:GP} menjelaskan struktur kategori
Grothendieck, sedangkan kategori abelian yang hingga secara lokal dan dibahas dalam
\S\ref{sec:locally-finite-Abel-cat} diperlukan oleh
\S\ref{sec:Tannaka-coend}. Pokok bagian terakhir ialah
Proposisi~\ref{prop:subcat-union} dari O.\ Gabber, yang buktinya memerlukan
sejumlah kecermatan.

% segment-id: o014.li2.appendix1.yoneda-density.h002
\section{Kepadatan pembenaman Yoneda}\label{sec:Yoneda}
% segment-id: o014.li2.appendix1.yoneda-density.p005
Pertama-tama, kita tinjau kerangka teoretis pembenaman Yoneda. Untuk
kategori $\mathcal{C}$ sebarang, definisikan kategori-kategori funktor
% segment-id: o014.li2.appendix1.yoneda-density.q001
\[ \mathcal{C}^\wedge := \cate{Set}^{\mathcal{C}^{\opp}}, \quad \mathcal{C}^\vee := \left(\cate{Set}^{\opp}\right)^{\mathcal{C}^{\opp}} = \left( \cate{Set}^{\mathcal{C}} \right)^{\opp}. \]
Keduanya saling dual:
$(\mathcal{C}^\vee)^{\opp}=(\mathcal{C}^{\opp})^\wedge$. Relatif terhadap
semesta Grothendieck $\mathcal{U}$ yang telah dipilih,
$\mathcal{C}^\wedge$ dan $\mathcal{C}^\vee$ biasanya kategori besar,
kecuali $\mathcal{C}$ kategori kecil.
\index{persoalan ukuran@persoalan ukuran}

% segment-id: o014.li2.appendix1.yoneda-density.p006
Istilah \emph{pembenaman Yoneda} berarti funktor-funktor berikut:
% segment-id: o014.li2.appendix1.yoneda-density.q002
\begin{align*}
	h_{\mathcal{C}}: \mathcal{C} & \longrightarrow \mathcal{C}^\wedge & k_{\mathcal{C}}: \mathcal{C} & \longrightarrow \mathcal{C}^\vee \\
	S & \longmapsto \Hom_{\mathcal{C}}(\cdot, S), & S & \longmapsto \Hom_{\mathcal{C}}(S, \cdot).
\end{align*}
\index{pembenaman Yoneda@pembenaman Yoneda (Yoneda embedding)}

% segment-id: o014.li2.appendix1.yoneda-density.p007
Berikut adalah pernyataan ulang \cite[Teorema 2.5.1]{Li1}.

% segment-id: o014.li2.appendix1.yoneda-density.th001
\begin{theorem}[Nobuo Yoneda]\label{prop:Yoneda}
	Untuk setiap $S\in\Obj(\mathcal{C})$,
	$A\in\Obj(\mathcal{C}^\wedge)$, dan
	$B\in\Obj(\mathcal{C}^\vee)$, pemetaan kanonik
	% segment-id: o014.li2.appendix1.yoneda-density.q003
	\[\begin{tikzcd}[row sep=tiny]
		\Hom_{\mathcal{C}^\wedge}\left(h_{\mathcal{C}}(S), A \right) \arrow[r] & A(S) \\
		{\left[ \Hom_{\mathcal{C}}(\cdot, S) \xrightarrow{\phi} A(\cdot) \right]} \arrow[mapsto, r] & \phi_S(\identity_S)
	\end{tikzcd} \quad \begin{tikzcd}[row sep=tiny]
		\Hom_{\mathcal{C}^\vee}\left(B, k_{\mathcal{C}}(S) \right) \arrow[equal, d] & \\
		\Hom_{\cate{Set}^{\mathcal{C}}}\left(k_{\mathcal{C}}(S), B \right) \arrow[r] & B(S) \\
		{\left[ \Hom_{\mathcal{C}}(S, \cdot) \xrightarrow{\psi} B(\cdot) \right]} \arrow[mapsto, r] & \psi_S(\identity_S)
	\end{tikzcd}\]
	semuanya bijeksi. Sebagai akibatnya, $h_{\mathcal{C}}$ dan
	$k_{\mathcal{C}}$ merupakan funktor setia penuh.
\end{theorem}

% segment-id: o014.li2.appendix1.yoneda-density.d001
\begin{definition}\label{def:representable-functor}
	\index{funktor!terwakili (representable)}
	Funktor $\mathcal{C}^{\opp}\to\cate{Set}$ (atau
	$\mathcal{C}\to\cate{Set}$) yang, hingga isomorfisme, berada dalam
	bayangan $h_{\mathcal{C}}$ (atau $k_{\mathcal{C}}$) disebut
	\emph{funktor terwakili}.
\end{definition}

% segment-id: o014.li2.appendix1.yoneda-density.p008
Walaupun $\mathcal{C}^\wedge$ (atau $\mathcal{C}^\vee$) jauh lebih besar
daripada $\mathcal{C}$, setiap objeknya dapat ditulis secara kanonik
sebagai kolimit (atau limit) funktor-funktor terwakili. Sebelum menyatakan
teorema kepadatan ini, kita memerlukan beberapa definisi.
\begin{itemize}
	\item Untuk setiap $A\in\Obj(\mathcal{C}^\wedge)$, definisikan kategori
	$(h_{\mathcal{C}}/A)$ seperti dalam Definisi~\ref{def:comma-category}.
	Objeknya berupa data
	$\underline S:=\bigl(S,h_{\mathcal{C}}(S)
	\xrightarrow{\phi_{\underline S}}A\bigr)$. Morfisme dari
	$\underline S$ ke $\underline S'$ ialah
	$f\in\Hom_{\mathcal{C}}(S,S')$ yang memenuhi
	$\phi_{\underline S}=\phi_{\underline S'}h_{\mathcal{C}}(f)$.

	\item Secara dual, untuk $B\in\Obj(\mathcal{C}^\vee)$ terdapat kategori
	$(B/k_{\mathcal{C}})$. Objeknya berupa data
	$\overline S:=\bigl(S,B\xrightarrow{\psi_{\overline S}}
	k_{\mathcal{C}}(S)\bigr)$, dengan $\psi_{\overline S}$ dipandang sebagai
	morfisme dalam $\mathcal{C}^\vee$; morfisme
	$\overline S\to\overline S'$ didefinisikan serupa.
\end{itemize}

% segment-id: o014.li2.appendix1.yoneda-density.th002
\begin{theorem}[Kepadatan]\label{prop:Yoneda-density}
	\index{pembenaman Yoneda!kepadatan (density)}
	Untuk setiap $A\in\Obj(\mathcal{C}^\wedge)$ dan
	$B\in\Obj(\mathcal{C}^\vee)$, keluarga morfisme
	$\phi_{\underline S}$ dan $\psi_{\overline S}$ di atas masing-masing
	memberikan isomorfisme kanonik dalam $\mathcal{C}^\wedge$ dan
	$\mathcal{C}^\vee$
	% segment-id: o014.li2.appendix1.yoneda-density.q004
	\[ \varinjlim_{\underline{S}} h_{\mathcal{C}}(S) \rightiso A, \quad B \rightiso \varprojlim_{\overline{S}} k_{\mathcal{C}}(S). \]
\end{theorem}
\begin{proof}
	Cukup buktikan pernyataan pertama. Untuk
	$A'\in\Obj(\mathcal{C}^\wedge)$ sebarang, terdapat pemetaan
	% segment-id: o014.li2.appendix1.yoneda-density.q005
	\begin{equation}\label{eqn:Yoneda-density-aux} \begin{aligned}
		\Hom_{\mathcal{C}^\wedge}(A, A') & \longrightarrow \left\{\begin{array}{l}
			\text{keluarga morfisme kompatibel}\; a'_{\underline{S}}: h_{\mathcal{C}}(S) \to A' , \\
			\text{dengan}\; \underline{S} = (S, \phi_{\underline{S}}) \in \Obj((h_{\mathcal{C}}/A))
		\end{array}\right\} \\
		\varphi & \longmapsto \left( a'_{\underline{S}} := \varphi \phi_{\underline{S}}\right)_{\underline{S}} .
	\end{aligned}\end{equation}

	Definisikan pemetaan baliknya sebagai berikut. Untuk data
	$(a'_{\underline S})_{\underline S}$, untuk setiap
	$S\in\Obj(\mathcal{C})$ dan $a_S\in A(S)$, ambil menurut
	Teorema~\ref{prop:Yoneda} morfisme yang bersesuaian
	$\phi:h_{\mathcal{C}}(S)\to A$, sehingga diperoleh
	$\underline S:=(S,\phi)\in\Obj((h_{\mathcal{C}}/A))$. Penetapan
	$a_S\mapsto a'_{\underline S}$ lalu memberikan pemetaan
	$A(S)\to A'(S)$. Klaimnya:
	\begin{compactitem}
		\item ketika $S$ bervariasi, pemetaan-pemetaan $A(S)\to A'(S)$
		memberikan morfisme $\varphi:A\to A'$ dalam $\mathcal{C}^\wedge$;
		\item pemetaan
		$(a'_{\underline S})_{\underline S}\mapsto\varphi$ dan pemetaan dalam
		\eqref{eqn:Yoneda-density-aux} saling invers.
	\end{compactitem}
	Seperti banyak hasil mengenai pembenaman Yoneda, verifikasi rinci hampir
	tautologis dan ditiadakan.

	Jadi \eqref{eqn:Yoneda-density-aux} bijeksi. Dalam kasus khusus $A'=A$,
	pemetaan tersebut membawa $\identity_A$ ke
	$(a'_{\underline S}=\phi_{\underline S})_{\underline S}$. Ini tepat
	memverifikasi bahwa $A$, bersama morfisme
	$\phi_{\underline S}:h_{\mathcal{C}}(S)\to A$, memenuhi sifat universal
	kolimit.
\end{proof}

% segment-id: o014.li2.appendix1.yoneda-density.p009
Pembaca juga dapat melihat pendekatan \cite[hlm. 76--77]{ML98}.

% segment-id: o014.li2.appendix1.yoneda-density.r001
\begin{remark}\label{rem:Yoneda-density}
	Isomorfisme mengenai $B$ dapat pula dinyatakan dalam
	$\cate{Set}^{\mathcal{C}}$ sebagai
	$\varinjlim_{\overline S}\Hom_{\mathcal{C}}(S,\mathord\cdot)
	\rightiso B$. Di sini $\overline S$ menjelajahi data
	$(S,\psi'_{\overline S})$ dengan $S\in\Obj(\mathcal{C})$ dan
	$\psi'_{\overline S}:\Hom_{\mathcal{C}}(S,\mathord\cdot)\to
	B(\mathord\cdot)$ morfisme dalam $\cate{Set}^{\mathcal{C}}$.
\end{remark}

% segment-id: o014.li2.appendix1.yoneda-density.p010
Salah satu penerapan Teorema~\ref{prop:Yoneda-density} ialah memperluas
funktor yang terdefinisi pada kategori kecil. Kita hanya menyatakan versi
$\mathcal{C}^\wedge$. Misalkan $\mathcal{C}$ kategori kecil,
$\mathcal{D}$ kategori kokomplet, yang boleh ``besar'' asalkan memiliki
semua kolimit kecil, dan $F:\mathcal{C}\to\mathcal{D}$ funktor. Definisikan
% segment-id: o014.li2.appendix1.yoneda-density.q006
\begin{equation}\label{eqn:Yoneda-F-tilde}\begin{aligned}
	\tilde{F}: \mathcal{C}^\wedge & \to \mathcal{D} \\
	X & \mapsto \varinjlim_{\underline{S}} F(S),
\end{aligned}\end{equation}
dengan
$\underline S=(S,\phi_{\underline S}:h_{\mathcal{C}}(S)\to X)$
menjelajahi kategori kecil $(h_{\mathcal{C}}/X)$. Bila $X$ berasal dari
objek $T$ dalam $\mathcal{C}$, kategori $(h_{\mathcal{C}}/X)$ memiliki
objek terminal $(T,h_{\mathcal{C}}(T)\rightiso X)$. Jadi makna perluasan
tersebut ialah komutativitas, hingga isomorfisme kanonik, diagram
% segment-id: o014.li2.appendix1.yoneda-density.q007
\[\begin{tikzcd}
	\mathcal{C}^\wedge \arrow[r, "{\tilde{F}}"] & \mathcal{D}. \\
	\mathcal{C} \arrow[u, "{h_{\mathcal{C}}}"] \arrow[ru, "F"'] &
\end{tikzcd}\]

% segment-id: o014.li2.appendix1.yoneda-density.p011
Sekarang kita tunjukkan bahwa $\tilde F$ selalu memiliki adjoin kanan dan
karena itu mempertahankan semua kolimit kecil. Fakta ini akan dipakai dalam
\S\ref{sec:Indization}.

% segment-id: o014.li2.appendix1.yoneda-density.pr001
\begin{proposition}\label{prop:Yoneda-F-tilde}
	Misalkan $\mathcal{C}$ kategori kecil dan $\mathcal{D}$ kategori
	kokomplet, yang boleh ``besar''. Untuk funktor
	$F:\mathcal{C}\to\mathcal{D}$, ambil perluasan kanonik
	$\tilde F:\mathcal{C}^\wedge\to\mathcal{D}$ menurut
	\eqref{eqn:Yoneda-F-tilde}. Maka $\tilde F$ memiliki adjoin kanan
	% segment-id: o014.li2.appendix1.yoneda-density.q008
	\[ \iHom(F, \cdot): \mathcal{D} \to \mathcal{C}^\wedge , \]
	yang didefinisikan oleh:
	\begin{itemize}
		\item $\iHom(F,Y)(S)=\Hom_{\mathcal{D}}(FS,Y)$ untuk setiap
		$S\in\Obj(\mathcal{C})$ dan $Y\in\Obj(\mathcal{D})$;
		\item morfisme $f:Y\to Y'$ menginduksi morfisme
		$f_*:\Hom(F(\mathord\cdot),Y)\to
		\Hom(F(\mathord\cdot),Y')$ dalam $\mathcal{C}^\wedge$.
	\end{itemize}
	Sebagai akibatnya, $\tilde F$ mempertahankan semua kolimit kecil.
\end{proposition}
\begin{proof}
	Untuk $X\in\Obj(\mathcal{C}^\wedge)$ dan
	$Y\in\Obj(\mathcal{D})$, dari
	$X=\varinjlim_{\underline S}h_{\mathcal{C}}(S)$ dan sifat universal
	kolimit diperoleh bijeksi kanonik
	% segment-id: o014.li2.appendix1.yoneda-density.q009
	\begin{multline*}
		\Hom_{\mathcal{C}^\wedge}(X, \iHom(F, Y)) \simeq \varprojlim_{\underline{S}} \Hom_{\mathcal{C}^\wedge}(h_{\mathcal{C}}(S), \iHom(F, Y)) \simeq \varprojlim_{\underline{S}} \Hom_{\mathcal{D}}(FS, Y) \\
		\simeq \Hom_{\mathcal{D}}(\varinjlim_{\underline{S}} FS, Y ) = \Hom_{\mathcal{D}}(\tilde{F}X, Y).
	\end{multline*}
	Jadi $\iHom(F,\mathord\cdot)$ benar-benar adjoin kanan dari
	$\tilde F:\mathcal{C}^\wedge\to\mathcal{D}$.
\end{proof}
